\fancyhead[L]{\small Third Proof}
\SourceLeaf{066}
\begin{center}
\textbf{THEOREM ON THE FACTORIZATION}\\[8pt]
\textbf{OF POLYNOMIALS}\\[8pt]
\textbf{INTO REAL FACTORS.}\\[8pt]
THIRD PROOF.\\[8pt]
SUPPLEMENT TO THE PRECEDING PAPER.\\[12pt]
\rule{3cm}{0.3pt}
\end{center}
After the preceding paper had already been printed, further reflection on the same subject led to a new proof of the theorem. Like the preceding proof, it is purely analytic, but it rests on entirely different principles and seems far preferable to it in simplicity. These few pages shall therefore be devoted to this \textit{third} proof.

\ThirdArticle{1}
Let the following function of the indeterminate $x$ be given:
\[
X=x^m+Ax^{m-1}+Bx^{m-2}+Cx^{m-3}+\text{etc.}+Lx+M
\]
where the coefficients $A,B,C$ etc. are fixed real quantities. Let $r,\varphi$ be other indeterminates, and put
\[
\begin{aligned}
&r^m\cos m\varphi+Ar^{m-1}\cos(m-1)\varphi\\
&\qquad+Br^{m-2}\cos(m-2)\varphi\\
&\qquad+Cr^{m-3}\cos(m-3)\varphi+\text{etc.}+Lr\cos\varphi+M=t\\
&r^m\sin m\varphi+Ar^{m-1}\sin(m-1)\varphi\\
&\qquad+Br^{m-2}\sin(m-2)\varphi\\
&\qquad+Cr^{m-3}\sin(m-3)\varphi+\text{etc.}+Lr\sin\varphi=u
\end{aligned}
\]
\SourceLeaf{067}
\[
\begin{aligned}
&mr^m\cos m\varphi+(m-1)Ar^{m-1}\cos(m-1)\varphi\\
&\qquad+(m-2)Br^{m-2}\cos(m-2)\varphi\\
&\qquad+(m-3)Cr^{m-3}\cos(m-3)\varphi+\text{etc.}+Lr\cos\varphi=t'\\
&mr^m\sin m\varphi+(m-1)Ar^{m-1}\sin(m-1)\varphi\\
&\qquad+(m-2)Br^{m-2}\sin(m-2)\varphi\\
&\qquad+(m-3)Cr^{m-3}\sin(m-3)\varphi+\text{etc.}+Lr\sin\varphi=u'\\[6pt]
&mmr^m\cos m\varphi+(m-1)^2Ar^{m-1}\cos(m-1)\varphi\\
&\qquad+(m-2)^2Br^{m-2}\cos(m-2)\varphi\\
&\qquad+(m-3)^2Cr^{m-3}\cos(m-3)\varphi+\text{etc.}+Lr\cos\varphi=t''\\
&mmr^m\sin m\varphi+(m-1)^2Ar^{m-1}\sin(m-1)\varphi\\
&\qquad+(m-2)^2Br^{m-2}\sin(m-2)\varphi\\
&\qquad+(m-3)^2Cr^{m-3}\sin(m-3)\varphi+\text{etc.}+Lr\sin\varphi=u''
\end{aligned}
\]
\[
\frac{(tt+uu)(tt''+uu'')+(tu'-ut')^2-(tt'+uu')^2}{r(tt+uu)^2}=y
\]
The factor $r$ can evidently be cancelled from the denominator of the last formula, since $t',u',t'',u''$ are divisible by it. Finally, let $R$ be a fixed positive quantity, arbitrary except that it must exceed the largest of the quantities
\[
mA\BareRoot{2},\quad\BareRoot{(mB\BareRoot{2})},\quad\BareRoot[3]{(mC\BareRoot{2})},\quad\BareRoot[4]{(mD\BareRoot{2})}\quad\text{etc.}
\]
disregarding the signs of the quantities $A,B,C$ etc., that is, replacing any negative quantities among them by positive ones. With these preparations, I assert that $tt'+uu'$ certainly takes a positive value when we put $r=R$, whatever (real) value is assigned to $\varphi$.

\textit{Proof.} Put
\[
\begin{aligned}
&R^m\cos45^\circ+AR^{m-1}\cos(45^\circ+\varphi)\\
&\qquad+BR^{m-2}\cos(45^\circ+2\varphi)\\
&\quad+CR^{m-3}\cos(45^\circ+3\varphi)+\text{etc.}\\
&\qquad+LR\cos(45^\circ+(m-1)\varphi)+M\cos(45^\circ+m\varphi)=T\\
&R^m\sin45^\circ+AR^{m-1}\sin(45^\circ+\varphi)\\
&\qquad+BR^{m-2}\sin(45^\circ+2\varphi)\\
&\quad+CR^{m-3}\sin(45^\circ+3\varphi)+\text{etc.}\\
&\qquad+LR\sin(45^\circ+(m-1)\varphi)+M\sin(45^\circ+m\varphi)=U\\
&mR^m\cos45^\circ+(m-1)AR^{m-1}\cos(45^\circ+\varphi)\\
&\qquad+(m-2)BR^{m-2}\cos(45^\circ+2\varphi)\\
&\quad+(m-3)CR^{m-3}\cos(45^\circ+3\varphi)+\text{etc.}\\
&\qquad+LR\cos(45^\circ+(m-1)\varphi)=T'\\
&mR^m\sin45^\circ+(m-1)AR^{m-1}\sin(45^\circ+\varphi)\\
&\qquad+(m-2)BR^{m-2}\sin(45^\circ+2\varphi)\\
&\quad+(m-3)CR^{m-3}\sin(45^\circ+3\varphi)+\text{etc.}\\
&\qquad+LR\sin(45^\circ+(m-1)\varphi)=U'
\end{aligned}
\]
It is evident that

I. $T$ is composed of the parts
\SourceLeaf{068}
\[
\begin{aligned}
&\frac{R^{m-1}}{m\BareRoot{2}}[R+mA\BareRoot{2}\, .\,\cos(45^\circ+\varphi)]\\
&+\frac{R^{m-2}}{m\BareRoot{2}}[RR+mB\BareRoot{2}\, .\,\cos(45^\circ+2\varphi)]\\
&+\frac{R^{m-3}}{m\BareRoot{2}}[R^3+mC\BareRoot{2}\, .\,\cos(45^\circ+3\varphi)]\\
&+\frac{R^{m-4}}{m\BareRoot{2}}[R^4+mD\BareRoot{2}\, .\,\cos(45^\circ+4\varphi)]\\
&+\text{etc.}
\end{aligned}
\]
each of which is readily seen to be positive for any fixed real value of $\varphi$. Hence $T$ necessarily takes a positive value. Similarly, it is proved that $U,T',U'$ are also positive, so that $TT'+UU'$ is necessarily a positive quantity as well.

II. For $r=R$, the functions $t,u,t',u'$ become, respectively,
\[
\begin{gathered}
T\cos(45^\circ+m\varphi)+U\sin(45^\circ+m\varphi)\\
T\sin(45^\circ+m\varphi)-U\cos(45^\circ+m\varphi)\\
T'\cos(45^\circ+m\varphi)+U'\sin(45^\circ+m\varphi)\\
T'\sin(45^\circ+m\varphi)-U'\cos(45^\circ+m\varphi)
\end{gathered}
\]
as is readily proved by expansion. Hence the value of the function $tt'+uu'$, for $r=R$, is found to be $=TT'+UU'$, and is therefore a positive quantity. Q. E. D.

The same formulas also show that the value of the function $tt+uu$, for $r=R$, is $TT+UU$, and is therefore positive. We conclude that for no value of $r$ greater than each of $mA\BareRoot{2},\BareRoot{(mB\BareRoot{2})},\BareRoot[3]{(mC\BareRoot{2})}$ etc. can $t=0,u=0$ hold simultaneously.

\ThirdArticle{2}
\textsc{Theorem.} \textit{Within the limits $r=0$ and $r=R$, and $\varphi=0$ and $\varphi=360^\circ$, there certainly exist values of the indeterminates $r,\varphi$ for which $t=0$ and $u=0$ hold simultaneously.}

\textit{Proof.} Suppose the theorem is false. Evidently, the value of $tt+uu$ must then be a positive quantity for all values of the indeterminates within the stated limits, and the value of $y$ must therefore always be finite. Consider the double integral
\[
\iint y\,\mathrm dr\,\mathrm d\varphi
\]
extended from $r=0$ to $r=R$, and from $\varphi=0$ to $\varphi=360^\circ$, which therefore takes a finite, fully determined value. This value, which
\SourceLeaf{069}
we shall denote by $\Omega$, must be the same whether we integrate first with respect to $\varphi$ and then with respect to $r$, or in the reverse order. But, treating $r$ as constant, we have the \textit{indefinite integral}
\[
\int y\,\mathrm d\varphi=\frac{tu'-ut'}{r(tt+uu)}
\]
as is readily confirmed by differentiation with respect to $\varphi$. No constant is to be added if we suppose the integral to begin at $\varphi=0$, since for $\varphi=0$ we have $\frac{tu'-ut'}{r(tt+uu)}=0$. Thus, since $\frac{tu'-ut'}{r(tt+uu)}$ evidently also vanishes for $\varphi=360^\circ$, the integral $\int y\,\mathrm d\varphi$ from $\varphi=0$ to $\varphi=360^\circ$ is $=0$, while $r$ remains indeterminate. It follows that $\Omega=0$.

\Needspace{20\baselineskip}
Likewise, treating $\varphi$ as constant, we have the indefinite integral
\[
\int y\,\mathrm dr=\frac{tt'+uu'}{tt+uu}
\]
as is just as readily confirmed by differentiation with respect to $r$. Here too no constant is to be added when the integral begins at $r=0$. Thus, by what was proved in the preceding article, the integral extended from $r=0$ to $r=R$ is $=\frac{TT'+UU'}{TT+UU}$, and so, by the theorem of that article, is always a positive quantity for every real value of $\varphi$. Hence $\Omega$ too, that is, the value of the integral
\[
\int\frac{TT'+UU'}{TT+UU}\,\mathrm d\varphi
\]
from $\varphi=0$ to $\varphi=360^\circ$, is necessarily a positive quantity\HistoricalNote{13}. This is absurd, since we previously found the same quantity to be $=0$. The supposition therefore cannot stand, and the truth of the theorem is established.

\ThirdArticle{3}
The function $X$ becomes $t+u\BareRoot{-1}$ under the substitution $x=r(\cos\varphi+\sin\varphi\, .\,\BareRoot{-1})$, and $t-u\BareRoot{-1}$ under the substitution $x=r(\cos\varphi-\sin\varphi\, .\,\BareRoot{-1})$. Thus, if for fixed values of $r,\varphi$, say $r=g,\varphi=G$, we have $t=0,u=0$ simultaneously (such values were proved to exist in the preceding article), $X$ takes the value 0 under each substitution
\SourceLeaf{070}
\[
x=g(\cos G+\sin G\, .\,\BareRoot{-1}),\quad x=g(\cos G-\sin G\, .\,\BareRoot{-1})
\]
and is therefore identically divisible by
\[
x-g(\cos G+\sin G\, .\,\BareRoot{-1}),\quad\text{and also by}\quad x-g(\cos G-\sin G\, .\,\BareRoot{-1})
\]
Whenever neither $\sin G=0$ nor $g=0$ holds, these divisors are unequal, and $X$ is therefore also divisible by their product
\[
xx-2g\cos G\, .\,x+gg
\]
But whenever either $\sin G=0$, and hence $\cos G=\pm1$, or $g=0$ holds, those factors are identical, namely $=x\mp g$. Thus the function $X$ certainly contains a real divisor of the second or first degree. Since the same conclusion applies again to the quotient, $X$ can be completely resolved into such factors. Q.~E.~D.

\Needspace{25\baselineskip}
\ThirdArticle{4}
Although we have already fully accomplished the task before us, it will not be superfluous to add a few further remarks on the reasoning of Article 2. From the supposition that $t$ and $u$ do not vanish simultaneously for any values of the indeterminates $r,\varphi$ within the limits stated there, we arrived at an unavoidable contradiction, from which we concluded that the supposition was false. This contradiction must therefore disappear if there really are values of $r,\varphi$ for which $t$ and $u$ simultaneously become $=0$. To clarify this further, we observe that for such values $tt+uu=0$, so that $y$ itself becomes infinite. We may therefore no longer treat the double integral $\iint y\,\mathrm dr\,\mathrm d\varphi$ as a quantity with a definite value. Generally speaking, if $\xi,\eta,\zeta$ denote indeterminate coordinates of points in space, the integral $\iint y\,\mathrm dr\,\mathrm d\varphi$ represents the volume of the solid contained between the five planes whose equations are
\[
\xi=0,\quad\eta=0,\quad\zeta=0,\quad\xi=R,\quad\eta=360^\circ
\]
and the surface whose equation is $\zeta=y$, taking as negative those parts in which the coordinates $\zeta$ are negative. But it is tacitly understood here that the sixth surface is \textit{continuous}. When this condition fails as $y$ becomes infinite, it may well happen that this concept has no meaning. In such a case there can be no question of summing the integral $\iint y\,\mathrm dr\,\mathrm d\varphi$, and it is therefore no wonder that analytic operations applied by blind calculation to empty concepts lead to absurdities.

\SourceLeaf{071}
The integration $\int y\,\mathrm d\varphi=\frac{tu'-ut'}{r(tt+uu)}$ is a true integration, that is, a summation, only insofar as $y$ is finite everywhere between the limits over which it extends; it is absurd if $y$ becomes infinite anywhere between those limits. If we evaluate an integral such as $\int\eta\,\mathrm d\xi$, which generally represents the area between the axis of abscissas and the curve whose ordinate is $=\eta$ for the abscissa $\xi$, by the usual rules, without regard to continuity, we very often become entangled in contradictions. For example, putting $\eta=\frac1{\xi\xi}$, analysis gives the integral $=C-\frac1\xi$, which correctly determines the area as long as the curve remains continuous. But continuity is interrupted at $\xi=0$, and if someone inappropriately asks for the magnitude of the area from a negative abscissa to a positive one, the formula will give the absurd answer that it is negative. We shall discuss more fully on another occasion what these and similar paradoxes of analysis mean.

\Needspace{22\baselineskip}
One more observation may be added here. When questions that may become absurd in certain cases are posed \textit{without restriction}, analysis very often provides for itself by making the answer partly indeterminate. Thus, for the value of the integral $\iint y\,\mathrm dr\,\mathrm d\varphi$ extended from $r=e$ to $r=f$, and from $\varphi=E$ to $\varphi=F$, if the value of $\frac ut$
\[
\begin{aligned}
&\text{for }r=e,\ \varphi=E\quad\text{is denoted by}\quad\theta\\
&\qquad r=e,\ \varphi=F\quad .\quad .\quad .\quad .\quad .\quad .\quad .\quad\theta'\\
&\qquad r=f,\ \varphi=E\quad .\quad .\quad .\quad .\quad .\quad .\quad .\quad\theta''\\
&\qquad r=f,\ \varphi=F\quad .\quad .\quad .\quad .\quad .\quad .\quad .\quad\theta'''
\end{aligned}
\]
analytic operations readily give
\[
\mathrm{Arc.\ tang}\,\theta-\mathrm{Arc.\ tang}\,\theta'-\mathrm{Arc.\ tang}\,\theta''+\mathrm{Arc.\ tang}\,\theta'''
\]
In fact, the integral can have a definite value only when $y$ always remains finite within the stated limits. This value is certainly contained in the formula given, but is not yet completely determined by it, since Arc. tang. is a multivalued function. It will be necessary to decide separately, by other considerations (which are not difficult), precisely which values of the function are to be used in a given case. On the other hand, whenever $y$ becomes infinite anywhere within the stated limits, the question of the value of the integral $\iint y\,\mathrm dr\,\mathrm d\varphi$ is absurd. If, despite this, you insist on forcing an answer from analysis, it will give now one answer, now another, according to the different methods; nevertheless, every one of these answers will be contained in the general formula given above.
\begin{center}\rule{3cm}{0.3pt}\end{center}
